% !TEX root = sommaire.tex

\begin{remerciements}

\footnotetext{This work has been supported by the French government, through the France 2030 investment plan managed by the Agence Nationale de la Recherche, as part of the "UCA DS4H" project, reference ANR-17-EURE-0004.}

Mes premiers remerciements vont à Enrico, qui m'a accompagné tout au long de mon parcours doctoral.

Je me suis toujours senti encouragé à explorer de nouvelles idées, sachant que je pouvais compter sur son expertise dans le vaste monde de la recherche. Un simple « merci » ne saurait suffire à exprimer toute ma gratitude pour la patience, le soutien et le professionnalisme dont il a toujours fait preuve à mon égard.

Je souhaiterais également remercier Cinzia, qui m'a soutenu dès le début. C'est d'ailleurs elle qui m'a donné l'impulsion nécessaire pour me lancer dans cette aventure. Je la remercie également de m'avoir donné l'occasion de participer à ses activités d'enseignement à l'université et d'y collaborer. J'ai trouvé cette expérience particulièrement enrichissante et je la garderai toujours parmi mes plus beaux souvenirs.

J'adresse un remerciement tout particulier à mon équipe, à commencer par Yves. Sa grande disponibilité et sa capacité d'écoute ne lui ont jamais fait perdre cette aura de grand maître avec laquelle je l'ai toujours considéré, ni cette autorité naturelle qui le caractérise en tant que professeur. Il m'a apporté des pistes de réflexion extrêmement précieuses, qui m'ont permis de prendre du recul, de réfléchir et de progresser.

Comment ne pas évoquer Laurent, qui m'a aidé à clarifier la rédaction de mes publications scientifiques et qui a toujours été prêt à partager ses expériences et ses anecdotes. Avec Thomas et Tamara, il a également contribué à entretenir ma passion pour les échecs.

Avec Jean-Christophe, Davide et Maria, j'ai partagé de nombreux moments autour de la musique : à parcourir des chansons, à expérimenter des arrangements, à enchaîner les fausses notes, à rire et, finalement, à trouver ensemble la bonne harmonie. Monter sur scène à l'occasion de la Fête de la Musique d'Inria a été un moment merveilleux et incroyablement gratifiant, que je garderai toujours précieusement en mémoire.

Je tiens également à remercier toutes les personnes que j'ai eu le plaisir de rencontrer au cours de mon parcours de thèse : Christine, Benjamin, Martin, Étienne et Matthieu.

Et maintenant, place à ma partie italienne.

Un pensiero colmo di affetto e gratitudine va a
mia mamma, che mi è sempre stata accanto, sostenendomi e sopportandomi a 360°
come nessun altro avrebbe mai potuto fare. 
Grazie anche a tutta la mia famiglia per esserci sempre stata.

Vorrei dedicare un pensiero speciale a nonno Lino, che mi ha affiancato durante
tutti gli studi, dalle scuole elementari fino all'Università. È sempre stato
tanto orgoglioso di me. Quanto mi mancano le nostre telefonate! Insieme a mia
mamma, è stato il mio sostenitore numero 1, incoraggiandomi in ogni circostanza.
Spero che, ovunque si trovi, possa vedere quest'ultimo traguardo del mio
percorso accademico e sapere quanto il suo sostegno abbia significato per me.

Grazie di cuore a tutti.


\end{remerciements}
